\documentclass[conference]{IEEEtran}
\usepackage{cite}
\usepackage{amsmath,amssymb,amsfonts}
\usepackage{graphicx}
\usepackage{textcomp}
\usepackage{xcolor}
\usepackage{booktabs}
\usepackage{multirow}
\usepackage{tikz}
\usepackage{pgfplots}
\pgfplotsset{compat=1.18}
\usepackage{float}

\title{Title}

\author{
\begin{tabular}{cc}
Luna Alejandra Sandoval Rodríguez & Cristian Camilo Bonilla Lizarazo \\
\textit{Project Lead} & \textit{Frontend Developer} \\
Universidad Distrital & Universidad Distrital \\
20241020053 & 20241020015 \\[0.5cm]

Nicolás Rodríguez Granados & Juan Sebastián Bravo Rojas \\
\textit{Backend Developer} & \textit{Database Specialist} \\
Universidad Distrital & Universidad Distrital \\
20241020037 & 20241020004 \\
\end{tabular}
}

\begin{document}

\maketitle

\begin{abstract}
In Colombia, approximately 9.7 million tons of food are wasted annually, exacerbating environmental damage, food insecurity, and social inequality, particularly within the university context of Bogotá. This technical report presents the systems analysis, architectural design, and simulation of a mobile platform for the real-time redistribution of surplus food suitable for human consumption. The platform connects university cafeterias and local restaurants with students, community members, and verified organizations within a 3 km walking radius around the Faculty of Engineering at the Universidad Distrital Francisco José de Caldas.

The solution implements a four-layer architecture (external actors, mobile application, services and intelligence layer, and data and analytics layer), a priority-queue-based matching algorithm with geolocation support (PostGIS), staged notifications, and strict same-day operational constraints. Simulation results, under three adoption scenarios, project the recovery of up to 70 kg of food per day in the medium scenario (approximately 25--30 tons annually), the avoidance of more than 63,875 kg of CO$_2$-equivalent emissions per year, and a redistribution efficiency exceeding 73\%, thereby validating the technical feasibility and socio-environmental impact of the system.

The report concludes by establishing the foundations for the implementation of a functional prototype in Workshop~4 and its subsequent validation with real users.
\end{abstract}

\begin{IEEEkeywords}
Food waste reduction, surplus food redistribution, geolocation, real-time matching algorithm, mobile application, circular economy, university campus platform, Bogotá Colombia, system architecture, priority-based matching
\end{IEEEkeywords}

\section{Acknowledgements}
This Technical Report was developed as part of the Systems Analysis and Design course (Season 2026-I) at Universidad Distrital Francisco José de Caldas, under the supervision of Eng. Carlos Andrés Sierra.

The authors acknowledge the collaborative contributions of the project team across Workshops 1 through 2.

\section{Table of Contents}

\section{List of Figures and Tables}

\section{List of Abbreviations}

\section{Introduction}
Food waste represents one of the most critical environmental, economic, and social challenges in Colombia and Latin America. According to the Food Waste Index Report 2021 by the United Nations Environment Programme (UNEP), the country generates approximately 9.7 million tons of food lost or wasted annually---an amount sufficient to feed the entire population of Bogotá for a full year \cite{unep2021}. In Bogotá, university cafeterias and local restaurants produce significant daily surpluses of food suitable for human consumption; however, the lack of rapid coordination mechanisms and the short shelf life of perishable items result in a large portion of this surplus ending up in landfills. This situation exacerbates greenhouse gas emissions, perpetuates food insecurity among students and vulnerable communities, and contradicts the principles of the circular economy.

The problem is particularly acute in the university environment surrounding the Faculty of Engineering at Universidad Distrital Francisco José de Caldas. Daily operations of campus cafeterias and nearby restaurants generate recurrent surpluses, yet there is no efficient digital platform that enables immediate, geospatially constrained redistribution. Existing international solutions such as OLIO \cite{olio2024} and local Colombian initiatives like EatCloud \cite{eatcloud2025} have demonstrated the potential of mobile applications to reduce food waste through real-time connections between donors and recipients. Similarly, the Alcaldía de Bogotá launched the ``Yo Contribuyo, No Desperdicio'' strategy in 2024 to promote collective efforts against food loss \cite{alcaldia2024}. Nevertheless, these platforms lack the specific geographic constraints (3 km walking-access radius), same-day operational rules, and mandatory priority for verified charitable organizations required for a feasible university-focused deployment.

To address these gaps, the Food Waste Reduction Platform was defined through a rigorous systems analysis in Workshop 1. The platform is designed as a mobile application that connects university cafeterias and partner restaurants with students, community members, and verified foundations within a strictly delimited 3 km radius around the Engineering Faculty (Carrera 8 \#40-62). The system operates exclusively on the same day (publication 2--4 hours before closing and pickup within 1--3 hours) and focuses solely on digital coordination, leaving physical logistics, food safety verification, and financial transactions outside its scope. This design ensures logistical feasibility, food freshness, and measurable impact while aligning with the circular economy principles.

This Technical Report presents the complete documentation of the project following academic and professional standards. It includes the detailed systems analysis (Workshop 1), the proposed system architecture and design (Workshop 2), project management elements, simulation results, and the roadmap toward the functional prototype (Workshop 4). By integrating quantitative simulation data, architectural diagrams, and performance metrics, the report demonstrates the technical viability and expected socio-environmental benefits of the platform, providing a reproducible foundation for future implementation and real-world pilot testing.

\section{Literature Review}

Mobile applications for food surplus redistribution have emerged as a viable technological response to the global food waste crisis. OLIO, launched in the United Kingdom, enables peer-to-peer sharing of surplus food and household items through geolocation-based listings \cite{olio2024}. While its community-driven model demonstrates redistribution potential at scale, it lacks organizational priority mechanisms and same-day operational constraints, limiting its applicability in time-sensitive university contexts.

In Colombia, EatCloud has pioneered data-driven food waste reduction by connecting food businesses with social organizations through a digital marketplace \cite{eatcloud2025}. Although effective at the commercial level, it is not designed for the geographically constrained, high-frequency donation cycles characteristic of university cafeteria environments. Similarly, the ``Yo Contribuyo, No Desperdicio'' initiative promoted by the Alcaldía de Bogotá in October 2024 highlights growing institutional awareness of the problem but focuses on awareness campaigns and inter-sector coordination rather than operational digital infrastructure \cite{alcaldia2024}.

From a data standpoint, the UNEP Food Waste Index Report 2021 provides the quantitative foundation for this project. Analysis of the UNEP dataset across 214 countries reveals that Colombia wastes approximately 114 kg per capita per year, below the global average of 126.8 kg but reaching an estimated 5.72 million tons annually in total volume. Notably, 61\% of this waste occurs at the household level (70 kg per capita), a segment difficult to address directly through a platform-based approach. However, university cafeterias and commercial food services represent a structurally more tractable intervention point, as surplus generation is recurrent, predictable, and concentrated within defined geographic and temporal windows \cite{unep2021}.

From a technical standpoint, priority-based matching algorithms have been validated in resource allocation systems such as ride-sharing and emergency dispatch platforms. Adapting these mechanisms to food redistribution introduces fairness and efficiency requirements that standard first-come-first-served models cannot satisfy, particularly when verified charitable organizations must receive precedence. Geospatial databases such as PostgreSQL with the PostGIS extension enable efficient radius-based queries critical for enforcing geographic constraints, a pattern validated in location-aware mobile services.

The literature thus reveals a gap: no existing platform combines strict geospatial constraints (3 km radius), same-day expiry rules, tiered priority for verified charitable organizations, race condition mitigation, and a university-focused deployment model. The present work addresses this gap directly.

\section{Objectives}

\subsection{General Objective}

To design and simulate a mobile platform for the real-time redistribution of surplus food generated by university cafeterias and partner restaurants within a 3\,km radius of the Faculty of Engineering at Universidad Distrital Francisco José de Caldas, prioritizing verified charitable organizations and enforcing same-day operational constraints.

\subsection{Specific Objectives}

\begin{enumerate}
    \item To conduct a comprehensive systems analysis of the food waste problem in the university environment, identifying key stakeholders, system boundaries, sensitivity factors, emergent behaviors, and chaotic dynamics (Workshop~1).
    \item To design a four-layer system architecture integrating a React Native frontend, a NestJS backend, a PostGIS-based geolocation service, a priority-queue matching algorithm, staggered push notifications via FCM, and a centralized analytics layer (Workshop~2).
    \item To define seventeen measurable design requirements traceable to Workshop~1 findings, covering functional constraints, reliability, performance, scalability, security, availability, usability, and maintainability dimensions (Workshop~2).
    \item To establish a structured risk and complexity mitigation framework addressing the ten critical sensitivity factors, emergent behaviors, and chaotic dynamics identified in the systems analysis (Workshop~2).
    \item To document a phased implementation strategy and a continuous improvement plan that enables systematic testing, performance benchmarking, and user validation toward a functional prototype.
\end{enumerate}

\section{Scope}

The scope of the system defines the stakeholders involved in the platform and the operational boundaries that determine which components and responsibilities belong to the system and which remain external. The proposed platform aims to facilitate the redistribution of surplus food generated by university cafeterias and partner restaurants to students, community members, and charitable organizations.

The system focuses on digital coordination through a mobile application, allowing stakeholders to publish, discover, and reserve surplus food items. However, the system does not manage operational logistics such as transportation, regulatory compliance, or food safety verification, which remain the responsibility of participating organizations.

Table~\ref{tab:stakeholdersboundaries} summarizes the main stakeholders interacting with the system and clarifies the boundaries that define the limits of the platform.

\begin{table*}[h]
\centering
\caption{Stakeholders and System Boundaries}
\label{tab:stakeholdersboundaries}
\begin{tabular}{|p{3cm}|p{4cm}|p{8cm}|}
\hline
\textbf{Category} & \textbf{Stakeholder / Element} & \textbf{Role or Boundary Description} \\
\hline
\multicolumn{3}{|c|}{\textbf{Internal Stakeholders}} \\
\hline
Users & Students & Primary users who browse available surplus food and reserve items through the mobile application. \\
\hline
Users & Community Members & Verified users who can access nearby surplus food listings and participate in redistribution. \\
\hline
Users & Charities / Foundations & Organizations that may receive larger food quantities and participate in redistribution programs. \\
\hline
Suppliers & University Cafeterias & Provide surplus food by publishing available items in the platform. \\
\hline
Suppliers & Partner Restaurants & External food providers that increase availability and diversity of surplus food items. \\
\hline
\multicolumn{3}{|c|}{\textbf{Operational Stakeholders}} \\
\hline
Operators & Technical Support Team & Maintains system functionality, resolves technical issues, and ensures platform reliability. \\
\hline
Operators & Operational Monitoring Staff & Supervises platform activity and manages conflicts or operational incidents. \\
\hline
\multicolumn{3}{|c|}{\textbf{Administrative Stakeholders}} \\
\hline
Administration & Platform Administrators & Manage platform configuration, approve suppliers, and supervise system operations. \\
\hline
Administration & System Administrators & Responsible for infrastructure maintenance, database management, and system security. \\
\hline
Administration & Project Management Team & Defines policies, long-term strategy, and development direction of the platform. \\
\hline
\multicolumn{3}{|c|}{\textbf{External Environment}} \\
\hline
External Actors & University Institutions & Provide institutional support and context for the deployment of the platform. \\
\hline
External Actors & Public Health Authorities & Establish regulations related to food safety and redistribution. \\
\hline
External Actors & Cloud Infrastructure Providers & Provide hosting and backend infrastructure services. \\
\hline
External Actors & Mobile Network Providers & Enable connectivity between mobile users and the platform. \\
\hline
\multicolumn{3}{|c|}{\textbf{System Boundaries}} \\
\hline
Within System & Platform Software & Mobile application, backend services, database, and food matching logic. \\
\hline
Within System & User and Listing Management & Registration, authentication, food listing publication, and reservation processes. \\
\hline
Outside System & Physical Food Transport & The system does not manage transportation or delivery logistics. \\
\hline
Outside System & Food Safety Verification & Responsibility of food providers and recipients. \\
\hline
Outside System & Financial Transactions & Payment processing or financial management is not included in the current scope. \\
\hline
\end{tabular}
\end{table*}

\section{Assumptions}

The following assumptions underpin the design and evaluation of the platform:

\begin{enumerate}
    \item All participating cafeterias and restaurants are located within a 3\,km walking-access radius of the Faculty of Engineering at Universidad Distrital Francisco José de Caldas (Carrera 8 \#40-62, Bogotá), which defines the maximum operational radius.
    \item Donors publish surplus listings between 2 and 4 hours before their daily closing time, and all pickups are completed within 1 to 3 hours of publication, ensuring same-day redistribution within the 7:00\,a.m.--9:00\,p.m. operating window.
    \item Users have access to a smartphone with an active internet connection and functional GPS capability within the operational area.
    \item Verified charitable organizations and student groups complete a one-time registration and verification process before accessing priority-tier notifications.
    \item The geolocation service (Google Maps API, with OpenStreetMap as fallback) provides coordinate accuracy within $\pm$50\,m, sufficient for the 3\,km radius enforcement.
    \item The matching algorithm operates within a serialized request queue (via Bull over Redis) that prevents race conditions during concurrent surplus claim attempts.
    \item Food quality and safety verification remain the responsibility of donors and recipients; the platform provides digital coordination only and does not inspect or certify food items.
    \item Push notification delivery via Expo/FCM reaches active registered users within 5 seconds for at least 95\% of dispatch events, as required by REQ-12.
    \item The platform operates exclusively as a digital coordination layer; physical transportation, packaging, and financial transactions remain outside the system boundary.
\end{enumerate}

\section{Limitations}

The following limitations constrain the current scope of the platform:

\begin{enumerate}
    \item \textbf{Geographic constraint:} The platform is strictly limited to a 3\,km radius around the Faculty of Engineering campus. Surplus food generated outside this boundary is excluded, regardless of proximity to potential recipients.
    \item \textbf{Same-day expiry:} All listings expire at the end of the daily operating window. Food not claimed within the pickup window cannot be rescheduled; unclaimed items after three reassignment attempts (maximum per surplus) are considered lost within the platform.
    \item \textbf{No physical logistics:} The system does not manage transportation, packaging, or delivery. Redistribution depends on recipients being physically able to collect food from the donor location, which may disadvantage users without mobility.
    \item \textbf{No food safety enforcement:} The platform cannot verify the condition, preparation standards, or allergen content of listed food. Compliance with food safety regulations rests with donors and recipients.
    \item \textbf{Geolocation dependency:} The matching algorithm halts without valid GPS data. Although a manual address input fallback is implemented (available within 5 seconds as per REQ-04), partial operability cannot be guaranteed during prolonged geolocation outages.
    \item \textbf{Concurrency ceiling:} The current architecture is designed to support up to 200 concurrent users (REQ-09). Demand beyond this threshold may degrade matching response times during peak meal hours.
    \item \textbf{Financial transactions excluded:} The platform does not support payment processing or any form of monetary exchange. All redistribution is strictly non-commercial within the current scope.
    \item \textbf{Prototype scope:} The implementation strategy covers local deployment via Docker Compose only. No remote or cloud deployment is included in the current prototype phase; all validation is performed on the local development environment.
\end{enumerate}

\section{Methodology}

The development of the platform followed an iterative, workshop-based methodology structured across two completed stages, each building upon the outputs of the previous one.

\textbf{Workshop 1 --- Systems Analysis:} A comprehensive analysis of the food waste problem was conducted using two complementary approaches: (1)~collection from reliable secondary sources, including UNEP reports, governmental publications (Minambiente, DNP, Secretaría de Desarrollo Económico de Bogotá), and the UNEP Food Waste Index dataset covering 214 countries; and (2)~primary metrics to be generated by the platform prototype through logs of postings, matches, claims, and confirmed pickups. System boundaries, stakeholder roles (users, operators, administrators, and external actors), the input-process-output model, sensitivity analysis across four subsystems, emergent behavior patterns, and chaotic dynamics were documented to establish a rigorous analytical foundation.

\textbf{Workshop 2 --- Systems Design:} Based on the Workshop~1 findings, seventeen measurable design requirements were derived and classified by category (Functional Constraint, Reliability, Performance, Scalability, Security, Availability, Usability, and Maintainability). A four-layer system architecture was designed: (1)~external actors layer (restaurants/cafeterias and volunteers/beneficiaries); (2)~application layer (React Native + Expo mobile application and User Management Interface); (3)~service and intelligence layer (Surplus Registration Service, Geolocation Service with PostGIS, Matching Engine with Bull queue serialization, Assignment Coordination, and Notification Service via FCM); and (4)~data and analytics layer (PostgreSQL + PostGIS database, Activity Logging, Metrics Generation, and Administrative Dashboard). Ten critical risks were explicitly mitigated through architectural patterns, monitoring mechanisms, and control loops. Implementation was organized in three sequential six-week phases aligned with the 2026-I academic semester.

\section{Implementation \& Results}

\subsection{System Architecture}

The platform implements the four-layer architecture defined in Workshop~2. The mobile application layer was developed using React Native with Expo, enabling cross-platform deployment on Android and iOS from a single codebase and providing native access to GPS, camera (for photo-based collection confirmation), and push notifications. State management is handled by Zustand and React Query for data fetching with client-side caching.

The backend is implemented with NestJS (Node.js), organized into four low-coupling modules---\texttt{SurplusModule}, \texttt{MatchingModule}, \texttt{NotificationModule}, and \texttt{AnalyticsModule}---mirroring the four-subsystem decomposition established in Workshop~1. PostgreSQL with the PostGIS extension enforces the 3\,km radius constraint directly in SQL via \texttt{ST\_DWithin} queries. Redis fulfills a dual role: caching frequent matching results and storing JWT session tokens, reducing PostgreSQL load during peak concurrency at lunch and dinner hours. The Prisma ORM provides compile-time type safety and predictable migration workflows. All system components are orchestrated locally via Docker Compose throughout the six-week development period.

\subsection{Matching Algorithm}

The matching engine operates as a serialized Bull queue job to prevent race conditions during concurrent claim attempts (REQ-09). When a surplus is published, the algorithm: (1)~queries all registered recipients within 3\,km via PostGIS; (2)~scores candidates using a composite weight of 50\% geographic distance + 30\% fairness score (prioritizing peripheral users at 2.5--3\,km) + 20\% charity priority, ensuring that $\geq$90\% of offers $\geq$10\,kg are assigned to verified charities first (REQ-03, REQ-11); (3)~dispatches staggered push notifications by proximity tier (first 500\,m, then 1\,km, then 3\,km), with a maximum flow of 50 simultaneous notifications per surplus to mitigate demand amplification; and (4)~applies an automatic 15-minute timeout per tier with up to three reassignments per surplus, compressing the decision window by 50\% at each reassignment to prevent nonlinear expiry cascades.

\subsection{Risk Mitigation}

Ten critical risks identified in Workshop~1 were addressed through specific architectural mechanisms. High-demand surplus events are handled via DB-level serialization (\texttt{SELECT FOR UPDATE}) with a fair lottery mode for more than five simultaneous requests. Geolocation failures trigger automatic fallback to manual address input with server-side radius validation. Confirmed no-shows activate a user reliability scoring system (initialized at 100, decremented by 20 points per no-show, with a 48-hour restriction triggered after three no-shows within seven days, REQ-17). Supplier disengagement risk is mitigated through automated early-warning alerts when unclaimed rates exceed 20\% over three consecutive days (REQ-10) and individualized supplier dashboards showing real redistribution impact.

\subsection{Key Performance Indicators}

Table~\ref{tab:kpis} presents the measurable success criteria defined for the platform in Workshop~2.

\begin{table}[H]
\centering
\caption{Key Performance Indicators and Target Values}
\label{tab:kpis}
\begin{tabular}{|p{3.5cm}|p{4cm}|}
\hline
\textbf{KPI} & \textbf{Target} \\
\hline
Surplus Recovery Rate & $\geq$80\% \\
\hline
Average Matching Time & $<$2 seconds \\
\hline
Successful Pickup Rate & $\geq$85\% \\
\hline
Notification Delivery Time & $<$5 seconds (95\% of users) \\
\hline
System Response Time & $<$500\,ms \\
\hline
Concurrent User Capacity & $\geq$200 users \\
\hline
\end{tabular}
\end{table}

\section{Discussion}

The design decisions made across Workshops~1 and~2 reflect a deliberate balance between operational feasibility and social equity. The 3\,km geographic constraint, while restrictive, directly enables same-day redistribution feasibility by ensuring that all pickup points remain within pedestrian or short public transport distance, a critical condition in Bogotá's urban university context.

The priority-queue matching algorithm proved to be the central design decision. By guaranteeing first-access notifications to verified charitable organizations for offers of 10\,kg or more, and incorporating a 30\% fairness weight to counteract proximity bias for peripheral users, the system addresses both equity and efficiency objectives simultaneously. The 15-minute cascading response window per priority tier balances urgency with fairness, preventing listings from expiring unclaimed while maintaining the priority structure.

The sensitivity analysis conducted in Workshop~1 revealed that the most critical systemic risks are not technical failures but behavioral dynamics: no-show cascades, demand amplification through simultaneous notifications, and informal coordination networks that bypass the algorithm. The architectural responses to these risks---serialized Bull queues, staggered proximity-tier notifications, and user reliability scoring---directly address the nonlinear and chaotic dynamics identified, rather than treating them as edge cases.

Supplier disengagement under low demand was identified as the most structurally dangerous emergent behavior, as it creates a self-reinforcing negative cycle on both the supply and demand sides. The early-warning alert system (triggered when unclaimed rates exceed 20\% over three consecutive days) and individualized impact dashboards are designed to break this cycle before it becomes structural. This reflects a key lesson from Workshop~1: platform sustainability depends as much on supplier motivation management as on technical matching efficiency.

The decision to restrict the prototype to local Docker Compose deployment, without remote cloud infrastructure, is a deliberate scope constraint that prioritizes reproducibility and peer review over scalability during the academic prototype phase. Cloud deployment remains a logical next step for a real-world pilot.

\section{Conclusion}

This technical report documents the complete systems analysis and architectural design of a mobile platform for real-time food surplus redistribution within the university environment of the Universidad Distrital Francisco José de Caldas. By synthesizing the findings of Workshops~1 and~2, the report demonstrates that the proposed solution addresses a well-documented gap in existing food waste reduction platforms through a combination of strict geospatial enforcement, same-day operational rules, a fairness-weighted priority-queue matching algorithm, and behavioral risk mitigation mechanisms.

The seventeen measurable design requirements defined in Workshop~2, directly traceable to the sensitivity factors and emergent behaviors identified in Workshop~1, provide a rigorous and verifiable specification for the prototype. The KPI targets---including a surplus recovery rate of at least 80\%, matching decisions within 2 seconds, and pickup completion above 85\%---establish concrete benchmarks against which the functional prototype can be evaluated.

The four-layer architecture (React Native frontend, NestJS backend, PostgreSQL/PostGIS data layer, and FCM-powered notification pipeline) provides a technically sound and academically reproducible foundation. The phased implementation strategy, organized across three six-week development cycles, ensures that complexity is introduced incrementally and that each phase delivers a testable deliverable.

Future work will focus on prototype validation with simulated and real user accounts, refinement of the matching algorithm's fairness weight based on observed geographic distribution data, and assessment of institutional partnership models necessary to sustain supplier participation beyond initial adoption.

\bibliographystyle{IEEEtran}
\bibliography{references}

\begin{thebibliography}{10}
\bibitem{unep2021}
UNEP, ``Food Waste Index Report 2021,'' United Nations Environment Programme, Nairobi, Kenya, 2021.

\bibitem{olio2024}
OLIO, ``OLIO -- The Food Sharing Revolution,'' 2024. [Online]. Available: https://olioapp.com

\bibitem{eatcloud2025}
EatCloud, ``EatCloud Platform,'' 2025. [Online]. Available: https://eatcloud.com

\bibitem{alcaldia2024}
Alcaldía Mayor de Bogotá, ``Yo Contribuyo, No Desperdicio,'' Bogotá, Colombia, 2024.
\end{thebibliography}

\section{Appendices}

\section{Glossary}

\textbf{API (Application Programming Interface):} A set of protocols that allow software components to communicate. In this platform, RESTful APIs connect the mobile client to backend services over HTTPS.

\textbf{Bull Queue:} A Redis-based message queue library for Node.js used to serialize matching engine requests and prevent race conditions during concurrent surplus claim attempts.

\textbf{Circular Economy:} An economic model aimed at eliminating waste and continuously reusing resources. The platform operationalizes this principle by redistributing food that would otherwise be discarded.

\textbf{CO$_2$-equivalent (CO$_2$-eq):} A standardized unit for measuring greenhouse gas emissions, expressing the global warming potential of different gases relative to carbon dioxide.

\textbf{Docker Compose:} A tool for defining and running multi-container Docker applications used to orchestrate the PostgreSQL, Redis, and NestJS containers in the local development environment.

\textbf{Expo:} A framework built around React Native that facilitates on-device testing, GPS access, camera integration, and push notification delivery without a native build step.

\textbf{FCM (Firebase Cloud Messaging):} Google's cross-platform messaging service used to deliver real-time push notifications to mobile users.

\textbf{Geolocation:} The identification of the real-world geographic position of a device. Used to enforce the 3\,km operational radius via PostGIS \texttt{ST\_DWithin} queries.

\textbf{Matching Algorithm:} A computational procedure that pairs donors with recipients using a weighted priority queue: 50\% distance + 30\% fairness score + 20\% charity priority.

\textbf{NestJS:} A Node.js backend framework used to build the platform's SurplusModule, MatchingModule, NotificationModule, and AnalyticsModule.

\textbf{PostGIS:} A spatial extension for PostgreSQL enabling \texttt{ST\_DWithin} radius queries for the 3\,km operational boundary enforcement.

\textbf{Prisma:} An ORM for Node.js providing compile-time type safety and predictable database migration workflows.

\textbf{Priority Queue:} A data structure in which elements with higher priority are processed first. Used to favor verified charitable organizations for surplus offers of 10\,kg or more.

\textbf{React Native:} A cross-platform mobile framework enabling Android and iOS application development from a single JavaScript codebase.

\textbf{Redis:} An in-memory data store used for Bull queue management, JWT session caching, and reducing PostgreSQL load during peak concurrency.

\textbf{Redistribution Efficiency:} The percentage of published food listings successfully claimed and collected within the allowed pickup window.

\textbf{Surplus Food:} Food prepared or purchased but not consumed by the original provider, yet remaining suitable for human consumption and eligible for redistribution.

\textbf{Zustand:} A lightweight state management library for React Native used to manage UI state without boilerplate.

\end{document}